\chapter{Modelos comparados}
\label{cap:modelos}

Los cuatro notebooks ejecutables comparten la misma familia algorítmica ---\emph{gradient
boosted trees}--- y difieren en el tratamiento de datos, el hardware y el protocolo de
validación. Esta homogeneidad es intencional: hace atribuible cualquier diferencia de AUC a la
ingeniería de variables. \cref{tab:modelos} resume las configuraciones.

\begin{table}[ht]
  \centering
  \caption{Configuraciones comparadas. ``d9/d12'' denota profundidad máxima del árbol.}
  \label{tab:modelos}
  \small
  \begin{tabular}{lllll}
    \toprule
    & Baseline CPU & Baseline GPU & LightGBM (nroman) & XGB d12 (cdeotte) \\
    \midrule
    Librería        & XGBoost 3.4  & XGBoost 3.4  & LightGBM 4.x & XGBoost 3.4 \\
    Dispositivo     & CPU (16 hilos) & RTX 2060    & CPU (16 hilos) & RTX 2060 \\
    Profundidad     & 9            & 9            & 31 hojas     & 12 \\
    Árboles         & 500          & 500          & $\le$5\,000 (ES) & 2\,000--5\,000 (ES) \\
    Tasa de aprend. & 0.05         & 0.05         & 0.0069       & 0.02 \\
    Features        & 432 crudas   & 432 crudas   & 298 (RFE + FE) & $\sim$200 (V + UID) \\
    Categóricas     & etiqueta     & etiqueta     & etiqueta     & etiqueta + frecuencia \\
    Validación      & holdout 75/25 & holdout 75/25 & TSSplit $\times$3 & GroupKFold mes $\times$3 \\
    \bottomrule
  \end{tabular}
\end{table}

\section{Baseline XGBoost (inversion, xhlulu)}

El punto de referencia: unión de tablas, relleno de faltantes con $-999$, etiquetado entero de
categóricas y \texttt{XGBClassifier} de profundidad 9 con 500 árboles. La variante GPU delega el
histograma y la construcción de árboles a la RTX 2060. En XGBoost moderno la invocación es:
\begin{lstlisting}
xgb.XGBClassifier(n_estimators=500, max_depth=9, learning_rate=0.05,
                  subsample=0.9, colsample_bytree=0.9, missing=-999,
                  tree_method='hist', device='cuda')   # ex 'gpu_hist'
\end{lstlisting}

\section{LightGBM con selección de variables (nroman)}

Parte de 120 variables ``útiles'' obtenidas por eliminación recursiva de características (RFE)
en un notebook hermano, añade codificación de frecuencia, variables de calendario,
interacciones y el decimal del monto, y entrena LightGBM con hiperparámetros de optimización
bayesiana (\cref{tab:modelos}), validación \texttt{TimeSeriesSplit} y parada temprana. El modelo
final re-entrena con todo el período y el número óptimo de iteraciones.

\section{XGBoost con ingeniería de entidades (cdeotte XGB\_95 / XGB\_96)}

Un único diseño de modelo (profundidad 12, 2\,000--5\,000 árboles, tasa 0.02, muestreo
0.8/0.4) evaluado en dos niveles de tratamiento de datos:

\begin{description}
  \item[XGB\_95 (sin UID):] selección de $\sim$120 columnas \texttt{V} por bloques de
        correlación, normalización de columnas \texttt{D}, codificación de frecuencia de
        tarjeta/dirección/correo, combinaciones \texttt{card1+addr1(+email)}, agregaciones
        \emph{por tarjeta} de monto y \texttt{D9}/\texttt{D11}, y descarte por consistencia
        temporal.
  \item[XGB\_96 (con magic UID):] todo lo anterior \emph{más} el pseudo-cliente de
        \cref{sec:uid} y sus $\sim$40 agregaciones. El modelo es idéntico al de XGB\_95:
        la ablación de \cref{sec:ablacion} aísla el aporte del UID.
\end{description}

El notebook original cierra con un post-procesado que suaviza las probabilidades promediando
dentro de cada UID (transacciones del mismo cliente deberían recibir scores similares); ese
paso requiere listas externas de limpieza de UIDs y queda fuera del alcance local, documentado
como trabajo futuro.

\section{Costo computacional}

La RTX 2060 (6\,GB) acelera el entrenamiento 4.4$\times$ respecto a 16 núcleos CPU en el
baseline, con AUC indistinguible: la GPU es una palanca de \emph{productividad} (más
iteraciones por hora), no de precisión. El límite de memoria también define fronteras
prácticas: el flujo \emph{RAPIDS cuDF} del mismo autor (todo el preprocesamiento en GPU,
15$\times$ más rápido) exige más de 6\,GB de VRAM y no fue ejecutable en este hardware;
queda documentado en el corpus descartado.
